\chapter{Cataplexy}

\section*{Definition}
Cataplexy is a sudden, bilateral episode of skeletal muscle tone loss triggered by strong emotions -- most commonly laughter, surprise, or anger -- with preserved consciousness and full recovery within seconds to minutes. It is pathognomonic for narcolepsy type~1 and is distinct from syncope, seizure, or psychogenic collapse.

\section*{Pathophysiology}
Cataplexy reflects an inappropriate intrusion of rapid eye movement (REM) sleep atonia into wakefulness. In people with narcolepsy type~1, autoimmune-mediated loss of hypothalamic orexin (hypocretin) neurons disrupts the boundary between wake and REM states. Orexin deficiency disinhibits brainstem regions that normally suppress muscle tone during REM -- including the locus coeruleus, dorsal raphe, and nucleus pontis oralis -- such that emotional stimuli gain aberrant access to the atonia-generating pathway, producing transient paralysis.

\section*{Causes}
\begin{itemize}
 \item Narcolepsy type~1 (formerly narcolepsy with cataplexy): Autoimmune destruction of orexin-producing neurons in the lateral hypothalamus, strongly associated with HLA-DQB1*06:02
 \item Narcolepsy type~2: Less common, typically without cataplexy; rare cases may later develop cataplexy with disease progression
 \item Secondary cataplexy: Hypothalamic lesions (tumors such as craniopharyngioma, germinoma; neurosurgery; stroke; demyelinating plaques of multiple sclerosis; traumatic brain injury; encephalitis, especially anti-Ma2 encephalitis)
 \item Niemann–Pick disease type~C: A lysosomal storage disorder that can present with gelastic cataplexy at young age
 \item Prader–Willi syndrome: Genetic disorder with reported associated cataplexy-like episodes
 \item Medication-induced: Withdrawal from anticataplectic agents (e.g., antidepressants, sodium oxybate) can cause rebound cataplexy; rare association with certain antipsychotics
\end{itemize}

\section*{When to See a Doctor}

\begin{itemize}
 \item Episodes of sudden muscle weakness with preserved awareness -- distinguish from syncope, seizure, or drop attack
 \item Cataplexy accompanied by excessive daytime sleepiness (Epworth Sleepiness Scale $>$ 10) to evaluate for narcolepsy
 \item New-onset cataplexy in the setting of known hypothalamic pathology, neurologic malignancy, or recent head trauma
 \item Cataplexy with hallucinations, sleep paralysis, or disrupted nocturnal sleep (the narcolepsy tetrad)
 \item Cataplexy beginning after initiation or withdrawal of psychotropic medication
\end{itemize}

\section*{Related Symptoms}
\begin{itemize}
 \item Excessive daytime sleepiness
 \item Hypnagogic hallucinations (visual, auditory, or tactile at sleep onset)
 \item Hypnopompic hallucinations (upon awakening)
 \item Sleep paralysis
 \item Fragmented nocturnal sleep
 \item Automatic behaviors (performing tasks without conscious awareness)
 \item Memory lapses
 \item Microsleep episodes
\end{itemize}
\textbf{Important:} Cataplexy is highly specific for narcolepsy type~1, but its absence does not exclude the diagnosis; about 10\% of narcolepsy people never develop cataplexy, and others may present first with excessive daytime sleepiness alone.

\section*{Self-Care / Home Management}
\begin{itemize}
 \item Identify and anticipate emotional triggers (laughter, surprise, anger, excitement) and sit or lean against a support at the first sign of an episode to prevent falls.
 \item Educate family, friends, and coworkers about the benign nature of episodes so they avoid unnecessary emergency response.
 \item Maintain a consistent sleep-wake schedule with strategic 15--20 minute naps to reduce overall sleepiness and cataplexy frequency.
 \item Avoid alcohol, sleep deprivation, and shift work which lower the threshold for cataplexy episodes.
\end{itemize}

\section*{How Your Doctor Will Evaluate You}
\begin{itemize}
 \item Detailed history of trigger events (always emotional), pattern of muscle involvement (knees, face, neck, or full body), preserved awareness during attacks, and complete recovery between episodes.
 \item Multiple sleep latency test (MSLT) demonstrating mean sleep latency $\le$ 8 minutes and $\ge$ 2 sleep-onset REM periods (SOREMPs) after nocturnal polysomnography (NPSG) ruling out other sleep disorders.
 \item Measurement of cerebrospinal fluid orexin-A (hypocretin-1) $\le$ 110~pg/mL is diagnostic for narcolepsy type~1 when MSLT is equivocal.
 \item HLA typing for DQB1*06:02 (strongly associated but not diagnostic alone -- 12--25\% of general population carries the allele).
\end{itemize}

\section*{Treatment Options}
\begin{itemize}
 \item Sodium oxybate (gamma-hydroxybutyrate, Xyrem/Xywav) is first-line therapy for cataplexy, consolidates sleep and reduces episode frequency.
 \item Antidepressants with anti-cataplectic effect: Venlafaxine (SNRI, most commonly used), fluoxetine (SSRI), clomipramine (tricyclic) -- these suppress REM sleep and reduce cataplexy.
 \item Pitolisant (histamine H3-receptor inverse agonist) improves wakefulness and reduces cataplexy as a non-controlled alternative.
 \item Stimulants (modafinil, armodafinil, methylphenidate) manage daytime sleepiness but do not directly reduce cataplexy; they are often combined with an anticataplectic agent.
\end{itemize}

\clearpage
